\documentclass[11pt,a4paper]{article}
\input{../template/utm-lab}

\newcommand{\labnumber}{3}
\newcommand{\labtitle}{The Command Line: Shell, Pipes and Text Filters}
\newcommand{\studentname}{Daniel Nicorici}
\newcommand{\studentgroup}{FAF-243}

\begin{document}

\utmtitlepage
\tableofcontents
\newpage

\plainsection{Purpose of the work}

The purpose of this lab is to work fluently in bash: to move fast with completion and history,
to predict what the shell does to a command line before the program runs, to send output and
errors exactly where I want them, and to answer real questions about a data set with one line
of small filters joined by pipes, finally turning the best one-liners into a script.

The work was done on my own laptop running Arch Linux, in GNU bash 5.3.15. The data set for
Part 4 was generated with the loop from the assignment:

\begin{lstlisting}
$ mkdir -p ~/os-lab3 && cd ~/os-lab3
$ for i in $(seq 1 500); do
    echo "2026-10-05 10:$(printf %02d $((i % 60))) user$((i % 7)) $([ $((i % 5)) -eq 0 ] && echo FAIL || echo OK) /page$((i % 4))"
  done > access.log
$ head -3 access.log ; wc -l access.log
2026-10-05 10:01 user1 OK /page1
2026-10-05 10:02 user2 OK /page2
2026-10-05 10:03 user3 OK /page3
500 access.log
\end{lstlisting}

The check prints \code{500 access.log}, so the file is correct.

\section{Part 1. Move fast, find help}

Pressing Tab asks bash to complete the current word. \code{/usr/sh} has only one match, so the
first Tab completed it to \code{/usr/share/}; the next Tab only rang the bell because many
entries start with \code{/usr/share/}, and the one after that offered to list all 229 of them.
\code{cd /et} plus Tab completed to \code{cd /etc/} at once.

\begin{lstlisting}[caption={Tab completion}]
$ ls /usr/sh[Tab]
$ ls /usr/share/[Tab]
$ ls /usr/share/[Tab]
Display all 229 possibilities? (y or n) n
$ cd /et[Tab]
$ cd /etc/
$ pwd
/etc
\end{lstlisting}

\begin{lstlisting}[caption={Output of type and which}]
$ type cd ; type ls ; type echo ; which ls
cd is a shell builtin
ls is aliased to `ls --color=auto'
echo is a shell builtin
/usr/bin/ls
\end{lstlisting}

\begin{lstlisting}[caption={The directories in PATH}]
$ echo $PATH
/home/Nico/.local/bin:/home/Nico/.local/bin:/home/Nico/.bun/bin:/usr/local/sbin:/usr/local/bin:/usr/bin:/opt/cuda/bin:/usr/lib/jvm/default/bin:/usr/bin/site_perl:/usr/bin/vendor_perl:/usr/bin/core_perl
$ echo $PATH | tr ':' '\n'
/home/Nico/.local/bin
/home/Nico/.local/bin
/home/Nico/.bun/bin
/usr/local/sbin
/usr/local/bin
/usr/bin
/opt/cuda/bin
/usr/lib/jvm/default/bin
/usr/bin/site_perl
/usr/bin/vendor_perl
/usr/bin/core_perl
$ echo $PATH | tr ':' '\n' | wc -l
11
\end{lstlisting}

\begin{lstlisting}[caption={Exit code of a missing command}]
$ nosuchprog ; echo "exit code: $?"
bash: nosuchprog: command not found
exit code: 127
\end{lstlisting}

After running three commands, Ctrl+R started a reverse search through the history: the first
letter \code{e} matched \code{date}, \code{ec} already matched \code{echo first command}, and
Enter ran that command again.

\begin{lstlisting}[caption={Reverse history search with Ctrl+R}]
$ echo first command
first command
$ date
Mon Oct  5 17:32:25 EEST 2026
$ whoami
Nico
(reverse-i-search)`ech': echo first command
first command
\end{lstlisting}

\begin{lstlisting}[caption={Finding help without a browser}]
$ ls --help | head -5
Usage: ls [OPTION]... [FILE]...
List information about the FILEs (the current directory by default).
Sort entries alphabetically if none of -cftuvSUX nor --sort is specified.

Mandatory arguments to long options are mandatory for short options too.
$ man -k "list directory"
dir (1)              - list directory contents
ls (1)               - list directory contents
vdir (1)             - list directory contents
$ help cd | head -2
cd: cd [-L|[-P [-e]]] [-@] [dir]
    Change the shell working directory.
$ man ls | grep -n '^ *-t,\? '
161:     -t     sort by time, newest first; see --time
\end{lstlisting}

\textbf{Observations:}
\begin{enumerate}
  \item \textbf{Built-ins.} \code{cd} and \code{echo} are shell built-ins; \code{ls} is an
  external program (\code{/usr/bin/ls}), which my \code{\textasciitilde/.bashrc} wraps in the alias
  \code{ls --color=auto}. \code{cd} \emph{must} be a built-in because the current directory
  belongs to each process: an external \code{cd} would run as a child process, change the
  child's directory and exit, leaving the shell exactly where it was.
  \item \textbf{PATH.} My PATH has 11 entries (Figure 1.3), but only 10 different
  directories: \code{\textasciitilde/.local/bin} is added twice, once by \code{\textasciitilde/.zshrc} and once by
  \code{\textasciitilde/.bashrc}. \code{ls} is found in \code{/usr/bin}.
  \item \textbf{Missing command.} A command that is not found returns exit code 127.
  \item \textbf{Sort by time.} \code{ls -t} sorts by modification time, newest first (line 161
  of \code{man ls}); \code{ls -lt} shows it in long format.
\end{enumerate}

\section{Part 2. What the shell does to your line}

\begin{lstlisting}[caption={Files created in the play directory}]
$ mkdir -p ~/os-lab3/play && cd ~/os-lab3/play
$ touch a.c b.c file1.txt file2.txt notes.txt "alfa beta.txt"
$ ls
 a.c  'alfa beta.txt'   b.c   file1.txt   file2.txt   notes.txt
\end{lstlisting}

\begin{lstlisting}[caption={Globbing results}]
$ echo *.c
a.c b.c
$ echo file?.txt
file1.txt file2.txt
$ echo [a-b]*
a.c alfa beta.txt b.c
$ echo *.xyz
*.xyz
\end{lstlisting}

\begin{lstlisting}[caption={The ls error on a name with a space, and the fix}]
$ ls alfa beta.txt
ls: cannot access 'alfa': No such file or directory
ls: cannot access 'beta.txt': No such file or directory
$ ls "alfa beta.txt"
'alfa beta.txt'
$ ls alfa\ beta.txt
'alfa beta.txt'
\end{lstlisting}

\begin{lstlisting}[caption={Single quotes, double quotes and backslash}]
$ echo '$HOME' "$HOME" \$HOME
$HOME /home/Nico $HOME
\end{lstlisting}

\begin{lstlisting}[caption={Brace expansion, arithmetic and command substitution}]
$ echo {1..5} ; echo $((7 * 6))
1 2 3 4 5
42
$ echo "today is $(date +%A), I am $(whoami)"
today is Monday, I am Nico
$ touch log{1..3}.txt && ls log*
log1.txt  log2.txt  log3.txt
\end{lstlisting}

\begin{table}[H]
  \centering
  \caption{Predictions compared with the real output}
  \footnotesize
  \begin{tabular}{@{}>{\raggedright\arraybackslash}p{0.26\textwidth}>{\raggedright\arraybackslash}p{0.19\textwidth}>{\raggedright\arraybackslash}p{0.19\textwidth}>{\raggedright\arraybackslash}p{0.28\textwidth}@{}}
    \toprule
    \textbf{Command} & \textbf{My prediction} & \textbf{Real output} & \textbf{Explanation} \\
    \midrule
    \code{echo *.c} & \code{a.c b.c} & \code{a.c b.c} & bash replaces the pattern with matching names \\
    \code{echo file?.txt} & \code{file1.txt file2.txt} & \code{file1.txt file2.txt} & \code{?} matches exactly one character \\
    \code{echo [a-b]*} & \code{a.c alfa beta.txt b.c} & \code{a.c alfa beta.txt b.c} & names starting with a or b, sorted \\
    \code{echo *.xyz} & \code{*.xyz} & \code{*.xyz} & no match, so bash leaves the pattern unchanged \\
    \code{ls alfa beta.txt} & two ``cannot access'' errors & two ``cannot access'' errors & the space splits the name into two arguments \\
    \code{ls "alfa beta.txt"} & \code{\textquotesingle{}alfa beta.txt\textquotesingle{}} & \code{\textquotesingle{}alfa beta.txt\textquotesingle{}} & quotes keep it one argument \\
    \code{echo \textquotesingle{}\$HOME\textquotesingle{} "\$HOME" \textbackslash\$HOME} & \code{\$HOME /home/Nico \$HOME} & \code{\$HOME /home/Nico \$HOME} & only double quotes allow expansion \\
    \code{echo \{1..5\} ; echo \$((7 * 6))} & \code{1 2 3 4 5}, \code{42} & \code{1 2 3 4 5}, \code{42} & brace expansion, then arithmetic \\
    \code{echo "today is \$(date +\%A), I am \$(whoami)"} & \code{today is Monday, I am Nico} & \code{today is Monday, I am Nico} & \code{\$(...)} is replaced by the command's output \\
    \code{touch log\{1..3\}.txt \&\& ls log*} & \code{log1.txt log2.txt log3.txt} & \code{log1.txt log2.txt log3.txt} & braces create three names, the glob finds them \\
    \bottomrule
  \end{tabular}
\end{table}

\textbf{Observations:}
\begin{enumerate}
  \item \textbf{Predictions.} Table 2.1 compares each prediction with the real output. The one
  that is easy to get wrong is \code{echo *.xyz}: when a pattern matches nothing, bash does not
  remove it but passes it through unchanged, so \code{echo} prints \code{*.xyz} literally.
  \item \textbf{Who produced the file names.} The shell. Before running \code{echo}, bash
  performs pathname expansion: it replaces \code{*.c} with the matching names in the current
  directory, so \code{echo} only receives \code{a.c} and \code{b.c} as ordinary arguments.
  \item \textbf{Why \code{ls alfa beta.txt} failed.} Bash splits the line into words at spaces,
  so \code{ls} received two arguments, \code{alfa} and \code{beta.txt}, and neither file
  exists. Two fixes: quote the name (\code{"alfa beta.txt"} or \code{\textquotesingle{}alfa beta.txt\textquotesingle{}}), or
  escape the space with a backslash (\code{alfa\textbackslash{} beta.txt}). A glob such as
  \code{alfa*} also works, because a glob result is never split again.
  \item \textbf{Single vs.\ double quotes.} Single quotes make everything inside literal, so
  \code{\textquotesingle{}\$HOME\textquotesingle{}} stays \code{\$HOME}. Double quotes keep the text as one word but still expand
  \code{\$VAR}, \code{\$(...)} and \code{\$((...))}, so \code{"\$HOME"} becomes
  \code{/home/Nico}. A backslash quotes only the next character.
\end{enumerate}

\code{ls} shows \code{\textquotesingle{}alfa beta.txt\textquotesingle{}} in quotes because, when its output goes to a terminal,
GNU \code{ls} quotes names that contain spaces or special characters, so that the name can be
copied and pasted back into the shell as one argument.

\section{Part 3. Redirection, chaining and jobs}

\begin{lstlisting}[caption={stdout and stderr in separate files}]
$ cd ~/os-lab3
$ ls /etc /nope > out.txt 2> err.txt
$ wc -l out.txt ; cat err.txt
176 out.txt
ls: cannot access '/nope': No such file or directory
\end{lstlisting}

\begin{lstlisting}[caption={Comparing the order of redirections}]
$ ls /etc /nope > all.txt 2>&1 ; tail -2 all.txt
xdg
zsh
$ ls /etc /nope > all.txt 2>&1
$ grep -c nope all.txt
1
$ ls /etc /nope 2>&1 > all.txt
ls: cannot access '/nope': No such file or directory
$ grep -c nope all.txt
0
\end{lstlisting}

\begin{lstlisting}[caption={The ; \&\& || operators}]
$ ls /nope 2> /dev/null || echo "failed with code $?"
failed with code 2
$ mkdir -p tmp && cd tmp && echo "I am in $(pwd)" ; cd ~/os-lab3
I am in /home/Nico/os-lab3/tmp
$ true && echo "&& ran: previous succeeded"
&& ran: previous succeeded
$ false && echo "this is skipped"
$ false || echo "|| ran: previous failed"
|| ran: previous failed
\end{lstlisting}

\begin{lstlisting}[caption={Running a background job and stopping it}]
$ sleep 100 &
[1] 130123
$ jobs ; kill %1 ; jobs
[1]+  Running                    sleep 100 &
[1]+  Running                    sleep 100 &
$ ls /etc /nope > all.txt 2>&1
[1]+  Terminated                 sleep 100
\end{lstlisting}

\begin{lstlisting}[caption={Appending and input redirection}]
$ echo first > notes.md ; echo second >> notes.md ; cat notes.md
first
second
$ wc -l < access.log
500
$ wc -l access.log
500 access.log
\end{lstlisting}

The second \code{jobs} in Figure 3.4 still shows \code{Running}: \code{kill} only sends the
signal, and the process had not died yet when \code{jobs} ran a moment later. Bash reported
\code{Terminated} just before the next prompt.

\textbf{Observations:}
\begin{enumerate}
  \item \textbf{Which file got what.} \code{err.txt} received the error message, because
  \code{2>} redirects file descriptor 2 (standard error). \code{out.txt} received the listing
  from standard output: 176 lines, the \code{/etc:} header plus 175 entries.
  \item \textbf{\code{> all.txt 2>\&1} vs.\ \code{2>\&1 > all.txt}.} Redirections are applied
  from left to right, and \code{2>\&1} means ``make stderr point where stdout points
  \emph{now}''. In the first form stdout is already the file, so both streams go into
  \code{all.txt} (\code{grep -c} finds the error once). In the second form stderr is copied
  from stdout while stdout is still the terminal, and only afterwards is stdout moved to the
  file, so the error appears on the screen and \code{all.txt} contains no error (count 0).
  \item \textbf{\code{wc -l < access.log}.} With \code{<} the shell opens the file and hands
  it to \code{wc} as standard input, so \code{wc} never learns a file name and prints only the
  number. With \code{wc -l access.log}, \code{wc} opens the file itself and prints its name.
  \item \textbf{Operators.} \code{;} runs the next command no matter what happened before.
  \code{\&\&} runs the next command only if the previous one succeeded (exit code 0).
  \code{||} runs the next command only if the previous one failed (non-zero exit code); here
  \code{ls} failed with code 2.
\end{enumerate}

\section{Part 4. Answer questions with one-liners}

Every answer is one command line, saved in \code{oneliners.sh} with the question as a comment
above it. Fields in \code{access.log} are separated by single spaces: 1 = date, 2 = time,
3 = user, 4 = result, 5 = page.

\begin{lstlisting}[caption={Number of failed requests}]
$ grep -c FAIL access.log
100
\end{lstlisting}

\begin{lstlisting}[caption={Different pages requested}]
$ cut -d' ' -f5 access.log | sort -u | nl
     1  /page0
     2  /page1
     3  /page2
     4  /page3
\end{lstlisting}

\begin{lstlisting}[caption={Requests per page}]
$ cut -d' ' -f5 access.log | sort | uniq -c
    125 /page0
    125 /page1
    125 /page2
    125 /page3
\end{lstlisting}

\begin{lstlisting}[caption={Users with the most failed requests}]
$ grep FAIL access.log | cut -d' ' -f3 | sort | uniq -c | sort -k1,1nr -k2,2 | awk 'NR == 1 { max = $1 } $1 == max'
     15 user3
     15 user5
\end{lstlisting}

\begin{lstlisting}[caption={Requests and failures of user3}]
$ echo "user3: $(grep -cw user3 access.log) requests, $(grep -w user3 access.log | grep -c FAIL) failed"
user3: 72 requests, 15 failed
\end{lstlisting}

\begin{lstlisting}[caption={The last 3 failed requests, time and user}]
$ grep FAIL access.log | tail -n 3 | cut -d' ' -f2,3
10:10 user0
10:15 user5
10:20 user3
\end{lstlisting}

\begin{lstlisting}[caption={Login shells in /etc/passwd}]
$ cut -d: -f7 /etc/passwd | sort | uniq -c | sort -rn
     24 /usr/bin/nologin
      2 /usr/bin/bash
      1 /usr/bin/zsh
      1 /usr/bin/git-shell
      1 /bin/false
\end{lstlisting}

\begin{lstlisting}[caption={ls and ls -l line counts}]
$ echo "ls: $(ls /etc | wc -l) lines, ls -l: $(ls -l /etc | wc -l) lines, first line of ls -l: $(ls -l /etc | head -n 1)"
ls: 175 lines, ls -l: 176 lines, first line of ls -l: total 1280
\end{lstlisting}

\begin{lstlisting}[caption={Full run of oneliners.sh (bonus lines 6--48 omitted)}]
$ bash oneliners.sh
100
     1  /page0
     2  /page1
     3  /page2
     4  /page3
    125 /page0
    125 /page1
    125 /page2
    125 /page3
     15 user3
     15 user5
user3: 72 requests, 15 failed
10:10 user0
10:15 user5
10:20 user3
     24 /usr/bin/nologin
      2 /usr/bin/bash
      1 /usr/bin/zsh
      1 /usr/bin/git-shell
      1 /bin/false
ls: 175 lines, ls -l: 176 lines, first line of ls -l: total 1280
     1  2026-10-05 10:10 user3 FAIL /page2
     2  2026-10-05 10:20 user6 FAIL /page0
     3  2026-10-05 10:30 user2 FAIL /page2
     4  2026-10-05 10:40 user5 FAIL /page0
     5  2026-10-05 10:50 user1 FAIL /page2
[... lines 6-48 omitted ...]
    49  2026-10-05 10:10 user0 FAIL /page2
    50  2026-10-05 10:20 user3 FAIL /page0
\end{lstlisting}

\begin{lstlisting}[caption={Bonus: failed requests at minutes ending in 0}]
$ grep -E '^[^ ]+ [0-9]{2}:[0-9]0 [^ ]+ FAIL ' access.log | nl | tail -n 2
    49  2026-10-05 10:10 user0 FAIL /page2
    50  2026-10-05 10:20 user3 FAIL /page0
\end{lstlisting}

\textbf{Answers:}
\begin{enumerate}[label=Q\arabic*.]
  \item 100 requests failed.
  \item 4 different pages: \code{/page0}, \code{/page1}, \code{/page2}, \code{/page3}.
  \code{sort -u} keeps one copy of each and \code{nl} numbers them, so the last number is the
  count.
  \item Each page got 125 requests; \(4 \times 125 = 500\), the number of lines in the log.
  \item There is a tie: \code{user3} and \code{user5} have 15 failed requests each. The final
  \code{awk} keeps every line whose count equals the first (largest) one, so a tie is shown
  instead of silently picking one user.
  \item \code{user3} made 72 requests, 15 of which failed. \code{-w} matches \code{user3} only
  as a whole word.
  \item The last three failures are at 10:10 (\code{user0}), 10:15 (\code{user5}) and 10:20
  (\code{user3}).
  \item Five login shells appear: \code{/usr/bin/nologin} (24 accounts, system users that may
  not log in), \code{/usr/bin/bash} (2), \code{/usr/bin/zsh} (1, my account),
  \code{/usr/bin/git-shell} (1) and \code{/bin/false} (1).
  \item My prediction was that \code{ls -l} adds one extra line, and it does: 175 vs.\ 176.
  When the output is a pipe, \code{ls} prints one name per line, while \code{ls -l} prints one
  line per entry \emph{plus} the header \code{total 1280}, the disk space used by the listed
  files in 1\,KiB blocks.
\end{enumerate}

\textbf{Bonus:} 50 lines match. The pattern \code{[0-9]\{2\}:[0-9]0} accepts any time whose
minute ends in 0, and \code{FAIL} must follow the user field. A minute ends in 0 when
\(i\) is a multiple of 10, and a request fails when \(i\) is a multiple of 5, so both hold for
the 50 multiples of 10 up to 500.

\textbf{Self-check:} the Q3 counts add up to 500, and the per-user failure counts behind Q4
(14 for \code{user0}, \code{user1}, \code{user2}, \code{user4}, \code{user6} and 15 for
\code{user3}, \code{user5}) add up to \(5 \times 14 + 2 \times 15 = 100\), the answer to Q1.
The file runs with \code{bash oneliners.sh} inside \code{\textasciitilde/os-lab3}.

\section{Part 5. Your first script}

The TODO is completed with a loop over \code{user0}--\code{user6}: brace expansion produces
the numbers 0 to 6, and for each user \code{grep -w} selects that user's lines and
\code{grep -c FAIL} counts the failures.

\begin{lstlisting}[caption={Contents of report.sh}]
#!/bin/bash
# usage: ./report.sh access.log
log=$1
if [ ! -f "$log" ]; then
	echo "no such file: $log" >&2
	exit 1
fi
total=$(wc -l < "$log")
fails=$(grep -c FAIL "$log")
echo "$total requests, $fails failed"
for n in {0..6}; do
	echo "user$n: $(grep -w "user$n" "$log" | grep -c FAIL) failures"
done
\end{lstlisting}

\begin{lstlisting}[caption={Running the script on access.log}]
$ chmod +x report.sh
$ ./report.sh access.log
500 requests, 100 failed
user0: 14 failures
user1: 14 failures
user2: 14 failures
user3: 15 failures
user4: 14 failures
user5: 15 failures
user6: 14 failures
\end{lstlisting}

\begin{lstlisting}[caption={Error handling for a missing file}]
$ ./report.sh nope.log ; echo "exit code: $?"
no such file: nope.log
exit code: 1
$ ./report.sh nope.log 2> /dev/null ; echo "exit code: $?"
exit code: 1
\end{lstlisting}

\begin{lstlisting}[caption={Running the script before chmod +x and without ./}]
$ ./report.sh access.log ; echo "exit code: $?"
bash: ./report.sh: Permission denied
exit code: 126
$ report.sh access.log
bash: report.sh: command not found
\end{lstlisting}

\textbf{Observations:}
\begin{enumerate}
  \item \textbf{The \code{\#!/bin/bash} line.} It is the ``shebang'': when the kernel is asked
  to execute the file, it sees \code{\#!} and runs \code{/bin/bash} with the script as its
  argument, so the line chooses the interpreter. Before \code{chmod +x} the file has no execute
  permission, so the kernel refuses to run it: \code{Permission denied}, exit code 126
  (Figure 5.4). It could still be run as \code{bash report.sh}.
  \item \textbf{Why \code{./report.sh}.} A command name without a slash is searched only in the
  PATH directories, and the current directory is not in PATH (Figure 1.3), so
  \code{report.sh} alone gives \code{command not found}. \code{./} is an explicit path to the
  file in the current directory.
  \item \textbf{\code{>\&2} and \code{exit 1}.} The error goes to standard error so that it
  never mixes with the real output: if the report is piped or saved to a file, the error still
  reaches the screen, and it can be discarded separately with \code{2> /dev/null}
  (Figure 5.3). \code{exit 1} returns a non-zero exit code, which tells the caller
  (\code{\&\&}, \code{||}, another script) that the script failed even when the message is
  hidden.
\end{enumerate}

\plainsection{Conclusions}

This lab showed how much the shell does before a program even starts. Bash splits the line into
words, expands globs, variables, braces, arithmetic and command substitutions, removes the
quotes, and sets up redirections, so the program only receives a ready list of arguments and
already-open file descriptors. Predicting each line before running it made these steps
visible, especially the cases that surprise, such as an unmatched glob that stays literal or a
space that splits one file name into two arguments.

Redirection is just rewiring file descriptors from left to right, which is why the order of
\code{>} and \code{2>\&1} matters, and exit codes are what \code{\&\&} and \code{||} decide on.
With these pieces, questions about 500 lines of log data were answered by short pipelines of
\code{grep}, \code{cut}, \code{sort}, \code{uniq} and \code{wc}, each tool doing one small
job. The same pipelines became a reusable script once it had a shebang, execute permission,
error messages on standard error and a meaningful exit code.

\end{document}
